\documentclass[../main.tex]{subfiles}

\begin{document}
\addcontentsline{toc}{chapter}{Preface}
\chapter*{Preface}

Linear algebra is a foundational branch of mathematics that
primarily investigates in linear equations, vectors, and their
relations. The application of linear algebra is immense that
numerous fields require foundational understanding of linear
algebra. For instance, it would not be possible to gain a
profound understanding in machine learning and data science
without the basics of linear algebra. This collection of notes
will discuss the essence of linear algebra in perspective of
teaching.


\addcontentsline{toc}{section}{About the Author}
\section*{\S\! About the Author}

\addcontentsline{toc}{section}{Background Information}
\section*{\S\! Background Information}

One of the most effective way to learn, at least for me, is to
teach. During the summer of 2026, I was enrolled in the course
XM511 from Stanford Pre-Collegiate University-Level Online Math,
instructed by Dr. Margarita Kanarsky. During the lectures,
I thought it would be a great idea to write notes on linear
algebra just as I wrote the notes for \href{null}{machine learning}.
Writing the teaching-style notes for machine learning allowed
me to differentiate what I truly know.

That being said, my notes do not represent the course. I have
included my understandings from numerous external sources and
restructured the notes itself to adhere to the
\href{https://www.stanford.edu/site/terms/}{Permission to Use
Materials}. Therefore, any errors in these notes are my
responsibility.


\addcontentsline{toc}{section}{Acknowledgement}
\section*{\S\! Acknowledgement}

\end{document}


